\documentclass[11pt]{letter}
\usepackage[margin=2.5cm]{geometry}
\usepackage{mathpazo}
\usepackage{url}
\usepackage{hyperref}
\hypersetup{colorlinks=true,urlcolor=blue,linkcolor=blue}

\date{May 28, 2026}

\signature{Yan Ma$^{1,2}$ \\ Lizhuo Zhang$^{3}$ \\ Xinjie Wu$^{4}$ \\ $^{1}$Changsha University of Science and Technology, Changsha 410114, China \\ $^{2}$Hunan Agricultural University, Changsha 410128, China \\ $^{3}$School of Information and Intelligence, Hunan Agricultural University \\ $^{4}$College of Liberal Arts and Social Sciences, City University of Hong Kong, Hong Kong, China \\ Corresponding author e-mail: zhanglizhuo@hunau.edu.cn}

\address{Hunan Agricultural University \\ Changsha 410128, China}

\begin{document}

\begin{letter}{Guest Editors \\ Special Issue: ``Applications Based on Symmetry in Adversarial Machine Learning'' \\ \textit{Symmetry} \\ MDPI}

\opening{Dear Guest Editors,}

We would like to submit the enclosed manuscript, ``Prompt Sensitivity as an Adversarial Vulnerability in CLIP-Family Models for Zero-Shot Classroom Behavior Analysis,'' for consideration as an original research article in the Special Issue on \textit{Applications Based on Symmetry in Adversarial Machine Learning} in \textit{Symmetry}.

\textbf{Background and motivation.}
Adversarial robustness research has predominantly focused on pixel-space perturbations of image inputs. However, vision--language foundation models such as CLIP introduce a second attack surface: the text prompt. In zero-shot recognition, users craft natural language descriptions of target categories without access to the test distribution. We demonstrate empirically that minor rephrasing of these descriptions---perturbations of a few words that carry no semantic intent to deceive---can silently and catastrophically degrade recognition performance. This form of \emph{text-domain adversarial-like perturbation} is practically relevant for any behavior analysis deployment built on foundation models, and it breaks the implicit assumption that a well-trained model will behave symmetrically under semantically equivalent inputs.

\textbf{Summary of contributions.}
This study makes three main contributions. First, using five CLIP-family ViT-Large backbones evaluated on a multi-scene classroom behavior benchmark (three public SCB subsets), we show that prompt-space perturbations can produce performance swings comparable to or exceeding backbone-level differences: on SigLIP2, an alternate-wording CAPE variant drops Hit@1 from 95.5\% to 31.4\% on the TeacherBehavior subset---a 64.1 percentage-point collapse triggered solely by text rewording. Second, we show that metric asymmetry compounds this vulnerability: the apparent zero-shot advantage over supervised baselines (favoring CAPE on lenient Hit@1) is reversed on class-discriminative multi-label metrics (Macro-F1, Sample-F1), revealing that standard reporting conventions can conceal rather than surface the model's fragility. Third, bootstrap confidence intervals and paired significance tests provide formal uncertainty quantification, showing that several widely cited headline gaps lack statistical support.

\textbf{Relevance to the Special Issue.}
The Special Issue calls for work on the robustness, security, and interpretability of modern learning systems under adversarial conditions, explicitly including behavior analysis and foundation models as target domains. Our work contributes a concrete empirical study of \emph{text-domain adversarial fragility} in vision--language foundation models applied to a real behavior analysis scenario, complementing existing pixel-space adversarial attack studies. The symmetry angle is direct: two prompts that are semantically near-equivalent to a human produce radically asymmetric model outputs, exposing a fundamental asymmetry in the model's prompt manifold that has direct implications for trustworthy deployment.

\textbf{Statements.}
This manuscript has not been published previously and is not under consideration by any other journal. All authors have read and approved the manuscript. The study reuses only existing public dataset releases and does not involve new human-subject data collection. The authors declare no conflicts of interest. This work was funded by the Education Department of Hunan Province, China (Project Number: HNJG-2022-0608).

The SCB datasets are publicly available at \url{https://huggingface.co/datasets/wintonYF/SCB-Dataset}. The reproducibility package (code, prompt templates, significance scripts) is available at \url{https://github.com/zhanglizhuo/scb5-zeroshot}.

We believe this manuscript will be of interest to readers working on adversarial machine learning, foundation model robustness, behavior analysis, and trustworthy AI. We appreciate your consideration.

\closing{Sincerely,}

\end{letter}
\end{document}